\section{Literature Review}


Multiple studies have reported substantial impact of the shale boom on local economies in the United States,
particularly in terms of employment and income(\cite{ihs2012america}; \cite{weinstein2011economic}).
Despite these economic benefits, however, other studies have raised concerns about the potential negative impacts of the shale boom on local education.
For example, \textcite{rickman2017shale} suggests that the influx of high-paying, low-skilled jobs in local labor market undermined educational outcomes based on a synthetic control method.
Among initial residents in West Virginia, Montana, and North Dakota, the study finds significant reductions in high school and college attainment among the states subject to the boom. 
Similarly, using county-level data on oil drilling activity and high school dropout rates in Texas, \textcite{carpenter2019texas} proposes that the shale gas boom in Texas led to increased dropout rates among immigrants, 
although they did not observe a statistically significant difference in the dropout rates for the overall population.

Our research is closest to \textcite{rickman2017shale} in terms of both research question and methodology. 
In fact, we also employ synthetic control methods to analyze the impact of the shale boom on educational outcomes.
However, we diverge from their research and many others
by (i) focusing on the heterogeneous effects of the shale boom across different household income quantiles, and (ii) mainly relying on staggered difference-in-differences (DiD) framework to analyze the impact of the shale boom on income and educational outcomes. -
To the best of our knowledge, no previous research has examined the heterogeneous effects of the shale boom on educational outcomes by employing newly developed staggered DiD methods, such as the ones proposed by \textcite{callaway2021difference} and \textcite{sun_estimating_2021}.
It is important to note that such methods can embrace multiple treatment groups and control groups
across different treatment periods, enabling a more comprehensive analysis of the shale boom's impact across the entire U.S. states.
